I cannot save `Q4P2.tex` in this read-only workspace. Here is the complete file content.

```latex
\documentclass{article}
\usepackage{pathfinder-note}

\title{Strategic Selection at a Calibrated Wall Gate}

\begin{document}
\maketitle

\section{The pair}

Q studies mechanisms for AI agents with private preferences and capabilities, allowing an agent to report a type and later deviate from its prescribed action. P gives a marginal split-conformal certificate: for one frozen wall fitter, a room-level Hausdorff bound makes every non-abstained decision at a wall-distance gate agree with the reference-wall gate \ledger{1, 3}.

\section{Connection pursued}

The question was whether P's certificate survives when a Q-style agent can select a report-contingent wall fitter and then choose where to act. Q places the report before the later private signal, while P calibrates only one frozen fitter; a report-selected fitter therefore needs its own analysis \ledger{2, 3, 9}.

\section{What was found}

\begin{claim}
Fix a finite menu of wall fitters \(F_j\) before calibration. Suppose complete calibration episodes and the new episode are exchangeable, and every fitter uses the same wall convention and reference frame. For episode \(i\), calibrate
\[
R_i=\max_j d_H(F_j(X_i),W_i^*).
\]
Let \(\varepsilon_\alpha\) be its split-conformal order statistic, using the infinity convention when the requested rank exceeds the calibration sample. With probability at least \(1-\alpha\), every menu estimate for the new episode is within \(\varepsilon_\alpha\) of the reference walls. On that one event, P's distance-to-set inequality certifies every non-abstained first-gate decision, regardless of which feasible report and scan point the agent selects.
\end{claim}

The argument applies P's exchangeable-score rank proof to the maximum over the fixed menu. For every \(j\) and point \(p\), the resulting coverage event gives
\[
\left|\operatorname{dist}(p,F_j(X_{\mathrm{new}}))
-\operatorname{dist}(p,W_{\mathrm{new}}^*)\right|
\leq\varepsilon_\alpha.
\]
Thus estimated distance plus \(\varepsilon_\alpha\) at most \(\delta\) certifies acceptance, while estimated distance minus \(\varepsilon_\alpha\) greater than \(\delta\) certifies rejection. One episode event covers all reports and points, so strategic selection does not require a pointwise union bound \ledger{5, 8, 14}.

If an independent wrapper restricts report \(j\) to a frozen finite point set \(C_j(X)\), the episode score can instead be
\[
T_i=\max_{j,\;p\in C_j(X_i)}
\left|\operatorname{dist}(p,F_j(X_i))
-\operatorname{dist}(p,W_i^*)\right|.
\]
The same rank argument covers every allowed report and point. Since \(T_i\leq R_i\) on each calibration episode, its conformal radius is no larger than the menu-Hausdorff radius computed from those episodes \ledger{16}. This sharper score depends on the wrapper enforcing the declared point sets.

\section{Objections and limits}

\begin{objection}
Separate marginal calibration of each fitter does not survive strategic fitter selection. If a room feature \(U\) is uniform on \(m\) values and only fitter \(j=U\) fails on that room, each fixed fitter fails with probability \(1/m\), while a selector that observes \(U\) chooses a failing fitter every time.
\end{objection}

The menu-wide maximum repairs this marginal selection problem. The counterexample's selector must know \(U\) before reporting to match Q's timing; allowing a later fitter choice is a stronger extension. P's statistical coverage event is also not Q's pathwise hard verification evidence \ledger{5, 7, 9}.

Coverage remains marginal across episodes. If deployment occurs on an event \(D\) of probability \(q>0\), the certificate alone gives only
\[
\Pr(\text{wrong}\mid D)\leq\min\{1,\alpha/q\}.
\]
Selection of rare failure rooms can make this bound sharp. In particular, menu-wide calibration does not establish low error conditional on deployment \ledger{4, 7, 8}.

Calibration also does not make an agent obey a selective fallback policy. For private signals \(g,b\), actions \(U\) for use and \(F\) for fallback, and intrinsic utility difference
\(\Delta(s)=U_t(U,s)-U_t(F,s)\), a state-independent action reward supports \(U\) on \(g\) and \(F\) on \(b\) only if
\[
\Delta(g)\geq r(F)-r(U)\geq\Delta(b).
\]
Hence weak implementation is possible exactly when
\(\Delta(g)\geq\Delta(b)\). This obstruction concerns the restricted reward instrument and a genuinely private signal; verifiable signals or enforceable signal-contingent permissions change the problem \ledger{6, 12, 14}.

\section{Outside sources and remaining work}

The prior-work check identified simultaneous conformal inference after data-dependent choice by Sarkar and Kuchibhotla, \url{https://arxiv.org/abs/2304.06158}; conformal prediction after model selection by Liang, Zhu, and Barber, \url{https://arxiv.org/abs/2408.07066}; and selection-conditional coverage by Bao and coauthors, \url{https://jmlr.org/papers/volume26/24-0452/24-0452.pdf}. These comparisons rule out presenting the generic maximum-score repair or selection warning as a standalone novelty claim \ledger{10, 12}.

The material is thin as a research result. Neither paper supplies ground-truthed wall episodes, a nonvacuous radius at P's \(\delta=0.15\,\mathrm m\), a complete fallback, or an implemented strategic report interface. The max-score certificate and two-action incentive test are corollaries; a substantive result would need a measured safety, abstention, and cost tradeoff under explicit incentives \ledger{8, 14, 16}.

The smallest next step is to freeze a small report-contingent fitter menu and its allowed scan-point sets, then collect disjoint ground-truthed calibration and test rooms. Compute both episode scores, their radii relative to \(\delta\), and the fraction of scan regions requiring fallback. That pilot would settle whether the proposed certificate has useful non-abstained coverage in the chosen population; incentive-compatible fallback would remain a separate question \ledger{5, 16}.

\end{document}
```